I am keen to apply for the 2026 Niels Bohr Quantum Summer School. The opportunity to meet and learn alongside an international cohort of researchers who are working on quantum algorithm design and software compilation is deeply appealing. This year’s focus connects directly to my PhD work on quantum cryptography, distributed quantum computing and quantum networks. My primary goal is to develop efficient, secure and verifiable protocols for distributed quantum computation that remain practical under realistic noise, restricted communication, and fault-tolerant implementation constraints. 

Some of my current and recent projects include designing and implementing protocols for fault-tolerant quantum Byzantine agreement, restricted quantum broadcasting through noisy channels, and efficient verifiable, blind, delegated quantum computing. I am interested in the interface between cryptographic protocol design and the software abstractions used to make protocols analyzable, compilable, and implementable. In that respect, I am especially eager to learn from Mingsheng Ying’s work on quantum program verification, Jens Palsberg’s perspective on quantum programming languages and compilation, and Daniel Gottesman’s work on quantum error correction and fault tolerance. These perspectives would help me sharpen the formal aspects of my own work, especially in translating cryptographic protocol descriptions into verifiable software-level abstractions.

Before moving into quantum cryptography, I was exposed to a number of components of the full quantum stack. Though I grew up in New Zealand, I was offered a full scholarship to complete my degree in the US where I am now a permanent resident. As an undergraduate I double majored in Mathematics and Physics while working as a research assistant doing experimental low temperature physics. In this role I synthesized graphene and carbon nanotubes through chemical vapor deposition and analyzed their properties using Raman spectroscopy. In addition to my lab and class time I competed as a track and field athlete and I was honored to be named graduation speaker.

Throughout my Master's degree at NYU I worked on the experimental side quantum computing with neutral atom ensembles. After my master's degree, I really wanted the opportunity to improve my programming skills which led to my position as a programmer/research assistant at Columbia University. Alongside pursuing a professional programming certificate through DataCamp, the research I was part of included building neural network classifiers for imputation using PyTorch, graph neural networks for causal inference, stochastic optimization algorithms, and Bayesian inference models. While at Columbia I participated in a mentorship program with the Quantum Open Source Foundation, where I co-authored a paper on Byzantine Agreement protocols utilizing quantum error correction that was accepted in the IEEE quantum networking 2025 conference. This experience convinced me to pursue my PhD and to focus on quantum cryptographic algorithms.

I would bring to the summer school a perspective grounded in quantum cryptography, quantum networks, and hybrid theory/software implementation. I am excited by the opportunity to learn from the lecturers, participate in the student challenges, and collaborate with an international cohort of researchers working on the foundations and implementation of quantum algorithms and software. Thank you for your consideration.

